\section*{Chapitre \ref{ch:dv:trading-volatility} --- Trading de la volatilité}

\begin{solution}{exo:dv:trading-volatility:1}
$767.5\times(0.02^2-0.18^2/252)=767.5\times0.000271=0.208$ pour la journée à 2~\%~; $-767.5\times0.18^2/252=-0.099$ pour la journée calme. Le mouvement quotidien
de seuil de rentabilité est $0.18/\sqrt{252}=1.13\%$.
\end{solution}

\begin{solution}{exo:dv:trading-volatility:2}
La variance réalisée est survenue lorsque le \omterm{def:dv:greeks-and-the-hedging-pnl:cashgamma}{cash gamma} du straddle était faible. Les semaines calmes (7.9~\% réalisée) se sont produites au strike, où le \omterm{def:dv:greeks-and-the-hedging-pnl:cashgamma}{cash gamma} était
d'environ 770 et le straddle payait du \omterm{def:dv:greeks-and-the-hedging-pnl:greeks}{theta} pour rien~: $-0.80$. La semaine agitée (45~\%) est survenue alors que le titre valait 107.8, où le \omterm{def:dv:greeks-and-the-hedging-pnl:cashgamma}{cash gamma} était de 39, soit
un vingtième du niveau initial~: $+0.02$.
\end{solution}

\begin{solution}{exo:dv:trading-volatility:3}
Carry $-1.88$~: la décroissance temporelle, à 21.79 inchangé, vaut $-1.59$ et le roll-down $-0.28$ (la volatilité baisse de 0.87 point jusqu'à 20.92 à deux mois). Sur une
courbe inversée, l'échéance la plus courte a la volatilité la plus élevée, donc une position longue remonte la courbe~: roll-down positif, qui compense une partie de la décroissance.
\end{solution}

\begin{solution}{exo:dv:trading-volatility:4}
Les straddles courts à un mois se décotent plus vite par unité de \omterm{def:dv:greeks-and-the-hedging-pnl:greeks}{vega} que le straddle long à trois mois, et la jambe à trois mois subit moins de roll-down par jour. Avec
la \omterm{def:dv:greeks-and-the-hedging-pnl:greeks}{vega} égalisée (1.73 straddles à un mois pour un straddle à trois mois), le \omterm{def:dv:greeks-and-the-hedging-pnl:greeks}{theta} net est positif~: 0.13 par jour. La contrepartie est un gamma court de $-0.17$, donc
un grand mouvement dans les prochains jours fait perdre de l'argent, comme un aplatissement de la courbe.
\end{solution}

\begin{solution}{exo:dv:trading-volatility:5}
$\sigma(K)=a+s\ln(K/S)/\sqrt\tau$. À $K$ et $a$ fixés, $\sigma_1-\sigma_0=s\bigl(\ln(K/S_1)-\ln(K/S_0)\bigr)/\sqrt\tau=-s\ln(S_1/S_0)/\sqrt\tau$, la même pour
chaque strike. Avec $s=-0.10$, $S_1/S_0=0.95$ et $\tau=0.25$~: $0.10\times\ln0.95/0.5=-0.0103$, soit une baisse d'environ un point.
\end{solution}

\begin{solution}{exo:dv:trading-volatility:6}
Le gamma et le \omterm{def:dv:greeks-and-the-hedging-pnl:greeks}{theta} dépendent de la trajectoire du sous-jacent et des cotations en début de journée, qui sont proches selon les deux règles lorsque le sous-jacent est proche
de sa référence. La \omterm{def:dv:greeks-and-the-hedging-pnl:greeks}{vega} et la \omterm{def:dv:greeks-and-the-hedging-pnl:greeks}{vanna} dépendent de $\delta\sigma_i$, la variation de la cotation de chaque option au cours de la journée. Sous \omterm{def:dv:implied-volatility-and-its-surface:sticky}{sticky strike}, elle ne dépend que du niveau,
alors que sous \omterm{def:dv:implied-volatility-and-its-surface:sticky}{sticky delta} elle dépend aussi du mouvement du sous-jacent, de sorte que les deux règles affectent le même mouvement à des colonnes différentes.
\end{solution}

\begin{solution}{exo:dv:trading-volatility:7}
Sans prime, la moyenne est de $-0.008\%$ du sous-jacent par mois, avec un écart-type de 1.33, soit une erreur type de $1.33/\sqrt{120}=0.12$~: indiscernable
de zéro. Le ratio de Sharpe vaut $-0.02$. La volatilité implicite vendue (21.5 en moyenne) dépasse la volatilité réalisée moyenne (20.6) parce qu'elle est la racine carrée de la
variance espérée, et, par l'inégalité de Jensen, cette racine dépasse l'espérance de la racine. En variance, les deux coïncident~: racines des moyennes quadratiques de 23.0 et 23.2. Les deux
points de prime représentent la totalité du $+0.44\%$.
\end{solution}

\begin{solution}{exo:dv:trading-volatility:8}
Le ratio de Sharpe et le taux de réussite décrivent le corps de la distribution, et non sa queue. Dans le programme du chapitre, les mêmes chiffres (1.15, 74~\%) coexistent avec un
pire mois de $-5.04\%$, 4.1 écarts-types sous la moyenne, et une asymétrie de $-1.35$. Dix années peuvent ne contenir aucun krach de la taille qui fixe le risque réel. Mesurer le
risque par scénario de stress (le 5 février 2018) et par des mesures de queue, non par la volatilité.
\end{solution}

\begin{solution}{pb:dv:trading-volatility:1}
\textbf{1.} 4.15~; que la variance réalisée pondérée par le gamma sur le mois dépasserait une volatilité de 18~\%.
\textbf{2.} $\pm1.57\%$ par jour~; $+3.28$.
\textbf{3.} Environ 770 ($\frac12\Gamma S^2$)~; un mouvement quotidien de $0.18/\sqrt{252}=1.13\%$ atteint le seuil de rentabilité.
\textbf{4.} 3.28 contre 3.05~: les termes d'ordre supérieur des mouvements quotidiens de 1.6~\%, que la somme d'ordre deux omet.
\textbf{5.} Le même résultat sur les trois trajectoires~: le notionnel multiplié par $25^2-18^2$ en points de variance, puisque le \omterm{def:dv:greeks-and-the-hedging-pnl:cashgamma}{cash gamma} d'un \omterm{def:dv:variance-swaps-and-volatility-derivatives:vs}{swap de variance} est constant.
\textbf{6.} Quinze jours de $+0.5\%$ et six de $\pm2.84\%$~: $15\times0.005^2+6\times0.0284^2=0.00521=0.25^2\times21/252$.
\textbf{7.} $-0.80$ pendant les semaines calmes, $+0.02$ pendant la semaine agitée~: $-0.78$.
\textbf{8.} Vers 107.8, soit 7.8~\% au-dessus du strike~: un \omterm{def:dv:greeks-and-the-hedging-pnl:cashgamma}{cash gamma} de 39 contre 770 à l'origine.
\textbf{9.} 2.87. Ses semaines calmes, au strike, versent du \omterm{def:dv:greeks-and-the-hedging-pnl:greeks}{theta} à plein gamma avec une réalisation de 7.9~\%~; la trajectoire régulière rapporte chaque jour.
\textbf{10.} L'\omterm{def:dv:trading-volatility:spread}{écart implicite-réalisé} n'est le P\&L espéré qu'en moyenne sur les trajectoires. Sur une trajectoire, le moment où la variance réalisée se produit par rapport au gamma décide~; un \omterm{def:dv:variance-swaps-and-volatility-derivatives:vs}{swap de variance}
isole cet écart.
\textbf{11.} Refixer le strike à la monnaie (vendre l'ancien straddle, en acheter un nouveau) aurait maintenu le gamma près de son niveau initial pendant la semaine agitée, au prix de
la fourchette et d'une nouvelle prime.
\textbf{12.} $-1.88$ sur le mois, dont $-0.28$ de roll-down.
\textbf{13.} 0.13 par jour, pour un gamma court de $-0.17$~: un grand mouvement ou un aplatissement de la courbe.
\textbf{14.} $-0.023$ en marquant en \omterm{def:dv:implied-volatility-and-its-surface:sticky}{sticky strike}, $+0.081$ en \omterm{def:dv:implied-volatility-and-its-surface:sticky}{sticky delta}.
\textbf{15.} L'explication par les \omterm{def:dv:greeks-and-the-hedging-pnl:greeks}{grecques} manque de 0.08 dans un sens comme dans l'autre~: les termes d'ordre trois d'un mouvement de 5~\% sont aussi grands que le P\&L de la journée.
\textbf{16.} Il ne s'agit pas du niveau~: l'action a bien réalisé 25. La transaction était un pari sur une quantité pondérée par la trajectoire, et la trajectoire a joué contre elle.
\textbf{17.} Le gamma a rapporté sur les semaines calmes moins que le \omterm{def:dv:greeks-and-the-hedging-pnl:greeks}{theta}, presque rien sur la semaine agitée, pas de \omterm{def:dv:greeks-and-the-hedging-pnl:greeks}{vega} (la volatilité implicite n'a pas bougé), et un petit résidu
non expliqué provenant des grands mouvements.
\textbf{18.} De nombreux petits gains et quelques pertes importantes (asymétrie de $-1.35$, pire mois de $-5.04\%$)~: la prime rémunère la prise en charge des krachs.
\textbf{19.} Avec la même réalisation de 25~\%, le straddle long perd 0.78 lorsque les mouvements surviennent les jours à faible gamma et gagne 3.28 lorsqu'ils sont répartis uniformément.
\textbf{20.} La variance réalisée pondérée par le \omterm{def:dv:greeks-and-the-hedging-pnl:cashgamma}{cash gamma} de l'option, moins la variance implicite sur les mêmes pondérations.
\end{solution}

\begin{solution}{iq:dv:trading-volatility:1}
À partir de l'écart entre les carrés des mouvements réalisés et la variance implicite, pondéré chaque jour par le \omterm{def:dv:greeks-and-the-hedging-pnl:cashgamma}{cash gamma} $\frac12\Gamma S^2$~: le gamma gagne lorsque les
mouvements sont grands, le \omterm{def:dv:greeks-and-the-hedging-pnl:greeks}{theta} perd lorsqu'ils sont petits~; s'y ajoutent la \omterm{def:dv:greeks-and-the-hedging-pnl:greeks}{vega} lorsque la volatilité implicite bouge, et les termes croisés.
\iqlookfor{l'écart pondéré par le \omterm{def:dv:greeks-and-the-hedging-pnl:cashgamma}{cash gamma}.}
\end{solution}

\begin{solution}{iq:dv:trading-volatility:2}
La variance réalisée est venue là où le gamma était faible~: le titre s'est éloigné du strike en silence puis a fortement bougé loin de lui, ou les grands mouvements sont survenus alors
que l'option avait beaucoup de temps devant elle et peu de gamma par unité de mouvement. Autres causes possibles~: la volatilité implicite a baissé (perte de \omterm{def:dv:greeks-and-the-hedging-pnl:greeks}{vega}), ou la couverture à intervalles discrets a manqué
les mouvements intrajournaliers.
\iqlookfor{dépendance à la trajectoire par le gamma, et la \omterm{def:dv:greeks-and-the-hedging-pnl:greeks}{vega}.}
\end{solution}

\begin{solution}{iq:dv:trading-volatility:3}
Long sur l'échéance la plus longue, short sur la plus courte, neutre en \omterm{def:dv:greeks-and-the-hedging-pnl:greeks}{vega}~: la décroissance plus rapide de la jambe courte donne un \omterm{def:dv:greeks-and-the-hedging-pnl:greeks}{theta} positif, et la jambe longue descend la courbe, ce qui coûte un
peu. Le carry net est en général positif~; le risque tient au gamma court et à un aplatissement ou une inversion de la courbe.
\iqlookfor{\omterm{def:dv:greeks-and-the-hedging-pnl:greeks}{theta}, roll-down, et le gamma court qui les finance.}
\end{solution}

\begin{solution}{iq:dv:trading-volatility:4}
Par option, à sa propre cotation~: delta, gamma, \omterm{def:dv:greeks-and-the-hedging-pnl:greeks}{theta}, \omterm{def:dv:greeks-and-the-hedging-pnl:greeks}{vega}, \omterm{def:dv:greeks-and-the-hedging-pnl:greeks}{vanna}, \omterm{def:dv:greeks-and-the-hedging-pnl:greeks}{volga} à partir des \omterm{def:dv:greeks-and-the-hedging-pnl:greeks}{grecques} de début de journée et des variations du sous-jacent et de chaque cotation au cours de la journée~; réévaluation complète
en parallèle. La part non expliquée est suivie~: petite et aléatoire, elle est acceptable, alors qu'une part grande ou persistante signale des facteurs de risque manquants, de grands mouvements, un changement de règle de valorisation ou
une erreur de saisie, chacune devant être examinée.
\iqlookfor{explication par option, réévaluation complète, part non expliquée comme contrôle.}
\end{solution}

\begin{solution}{iq:dv:trading-volatility:5}
Aucune des deux n'est juste en général~: ce sont deux hypothèses sur la façon dont la surface évolue avec le sous-jacent. Le P\&L total de la journée suit les cotations, et la répartition
suit la règle. Retenir la règle que le marché a suivie pour ce sous-jacent (chapitre~7) et signaler l'exposition à la \omterm{def:dv:greeks-and-the-hedging-pnl:greeks}{vanna} qu'elle crée comme un risque.
\iqlookfor{règle de valorisation comme hypothèse, et son effet sur la répartition.}
\end{solution}

\begin{solution}{iq:dv:trading-volatility:6}
Non~: il s'agit de la rémunération de la prise en charge du risque de krach et de la fourniture d'une assurance. Les rendements ont une longue queue gauche. Dimensionner par la perte en scénario de stress (un 5 février 2018 ou pire), et non par
la volatilité~; plafonner le gamma court et la \omterm{def:dv:greeks-and-the-hedging-pnl:greeks}{vega}~; privilégier les structures à risque défini~; garder du capital pour la queue.
\iqlookfor{rémunération du risque de queue, dimensionnement par scénario de stress.}
\end{solution}
